\documentclass{amsart}
% For dealing with references we use the comment environment
\usepackage{verbatim}
\newenvironment{reference}{\comment}{\endcomment}
%\newenvironment{reference}{}{}
\newenvironment{slogan}{\comment}{\endcomment}
\newenvironment{history}{\comment}{\endcomment}

% For commutative diagrams we use Xy-pic
\usepackage[all]{xy}

% We use 2cell for 2-commutative diagrams.
\xyoption{2cell}
\UseAllTwocells

% We use multicol for the list of chapters between chapters
\usepackage{multicol}

% This is generally recommended for better output
\usepackage{lmodern}
\usepackage[T1]{fontenc}

% For cross-file-references

% Package for hypertext links:
\usepackage{hyperref}

% For any local file, say "hello.tex" you want to link to please

% Theorem environments.
%
\theoremstyle{plain}
\newtheorem{theorem}[subsection]{Theorem}
\newtheorem{proposition}[subsection]{Proposition}
\newtheorem{lemma}[subsection]{Lemma}

\theoremstyle{definition}
\newtheorem{definition}[subsection]{Definition}
\newtheorem{example}[subsection]{Example}
\newtheorem{exercise}[subsection]{Exercise}
\newtheorem{situation}[subsection]{Situation}

\theoremstyle{remark}
\newtheorem{remark}[subsection]{Remark}
\newtheorem{remarks}[subsection]{Remarks}

\numberwithin{equation}{subsection}

% Macros
%
\def\lim{\mathop{\mathrm{lim}}\nolimits}
\def\colim{\mathop{\mathrm{colim}}\nolimits}
\def\Spec{\mathop{\mathrm{Spec}}}
\def\Hom{\mathop{\mathrm{Hom}}\nolimits}
\def\Ext{\mathop{\mathrm{Ext}}\nolimits}
\def\SheafHom{\mathop{\mathcal{H}\!\mathit{om}}\nolimits}
\def\SheafExt{\mathop{\mathcal{E}\!\mathit{xt}}\nolimits}
\def\Sch{\mathit{Sch}}
\def\Mor{\mathop{\mathrm{Mor}}\nolimits}
\def\Ob{\mathop{\mathrm{Ob}}\nolimits}
\def\Sh{\mathop{\mathit{Sh}}\nolimits}
\def\NL{\mathop{N\!L}\nolimits}
\def\CH{\mathop{\mathrm{CH}}\nolimits}
\def\proetale{{pro\text{-}\acute{e}tale}}
\def\etale{{\acute{e}tale}}
\def\QCoh{\mathit{QCoh}}
\def\Ker{\mathop{\mathrm{Ker}}}
\def\Im{\mathop{\mathrm{Im}}}
\def\Coker{\mathop{\mathrm{Coker}}}
\def\Coim{\mathop{\mathrm{Coim}}}

% Boxtimes
%
\DeclareMathSymbol{\boxtimes}{\mathbin}{AMSa}{"02}

%
% Macros for moduli stacks/spaces
%
\def\QCohstack{\mathcal{QC}\!\mathit{oh}}
\def\Cohstack{\mathcal{C}\!\mathit{oh}}
\def\Spacesstack{\mathcal{S}\!\mathit{paces}}
\def\Quotfunctor{\mathrm{Quot}}
\def\Hilbfunctor{\mathrm{Hilb}}
\def\Curvesstack{\mathcal{C}\!\mathit{urves}}
\def\Polarizedstack{\mathcal{P}\!\mathit{olarized}}
\def\Complexesstack{\mathcal{C}\!\mathit{omplexes}}
% \Pic is the operator that assigns to X its picard group, usage \Pic(X)
% \Picardstack_{X/B} denotes the Picard stack of X over B
% \Picardfunctor_{X/B} denotes the Picard functor of X over B
\def\Pic{\mathop{\mathrm{Pic}}\nolimits}
\def\Picardstack{\mathcal{P}\!\mathit{ic}}
\def\Picardfunctor{\mathrm{Pic}}
\def\Deformationcategory{\mathcal{D}\!\mathit{ef}}

\setlength{\emergencystretch}{3em}
\begin{document}
\title[Stacks Project, Tag 0GVB]{429 --- Prescribed sections with controlled codimension of degeneracy}
\maketitle
\noindent Copyright (C) 2005--2025 Johan de Jong. Stacks Project authors. Source: \href{https://stacks.math.columbia.edu/tag/0GVB}{Tag 0GVB}.

\noindent Editorial difficulty note (not author text). Difficulty H2: A double induction constructs auxiliary sections avoiding generic points of the existing bad components. Chinese-remainder interpolation combines them while preserving prescribed fibre values, and a component-by-component argument raises the codimension of every new degeneracy locus..

\noindent Editorial provenance note (not author text). History: excerpt from revision \texttt{a0444\allowbreak 6e57e\allowbreak c1fbc\allowbreak 252a8\allowbreak 71afc\allowbreak ec775\allowbreak 2fb28\allowbreak 07b14}, more-algebra.tex, lines 38808--38917. Reference presentation was adapted; original-author context and core support blocks were retained or added. The mathematical wording is unchanged. Licensed under GNU FDL 1.2 or later; no invariant sections or cover texts. See ../licenses/COPYING. Context blocks reproduce author definitions and supporting results; they are not additional counted problems.

\begin{situation}
\label{situation-splitting}
Here $R$ is a ring and $M$ is a finitely presented $R$-module.
Denote $\Omega \subset \Spec(R)$ the set of closed points
with the induced topology. For $x \in \Omega$ denote
$M(x) = M/xM$ the fibre of $M$ at $x$. This is a finite
dimensional vector space over the residue field $\kappa(x)$ at $x$.
Given $s \in M$ we denote $s(x)$ the image of $s$ in $M(x)$.
\end{situation}

\begin{proposition}
\label{proposition-splitting}
\begin{reference}
\href{https://stacks.math.columbia.edu/bibliography}{Bibliography: Serre-projective (Theorem 2)}
\end{reference}
In Situation \ref{situation-splitting} assume $\Omega$ is
a Noetherian topological space. Let $s_1, \ldots, s_h \in M$.
Let $Z(s_1, \ldots, s_h) \subset F \subset \Omega$ be closed.
Let $x_1, \ldots, x_n \in F$ be pairwise distinct.
Let $v_i  \in V(x_i)$.
Let $k \geq 0$ be an integer such that
$$
(*)\quad h + k \leq \dim_{\kappa(x)} V(x)\text{ for all }x \in \Omega
$$
Then there exist $s \in M$ and $F' \subset \Omega$ closed such that
\begin{enumerate}
\item[(a)] $s(x_i)$ maps to $v_i$,
\item[(b)] $Z(s_1, \ldots, s_h, s) \subset F \cup F'$, and
\item[(c)] every irreducible component of $F'$ has
codimension $\geq k$ in $\Omega$.
\end{enumerate}
\end{proposition}

\begin{proof}
We note that codimension was defined in
Topology, Section \href{https://stacks.math.columbia.edu/tag/02I0}{section catenary spaces (Tag 02I0)}
and that we will use some results on Noetherian topological
spaces contained in
Topology, Section \href{https://stacks.math.columbia.edu/tag/0050}{section noetherian (Tag 0050)}.

\medskip\noindent
The proof is by induction on $k$. If $k = 0$, then we choose
$s \in M$ as in Lemma \href{https://stacks.math.columbia.edu/tag/0GVA}{lemma choose values (Tag 0GVA)} and we choose
$F' = \Omega$.

\medskip\noindent
Assume $k > 0$. By our induction hypothesis we may
choose $u \in M$ and $G \subset \Omega$ closed satisfying
(a), (b), (c) for $s_1, \ldots, s_h$, $F$, $x_1, \ldots, x_n$,
$v_1, \ldots, v_n$, and $k - 1$.

\medskip\noindent
Let $G = G_1 \cup \ldots \cup G_m$ be the decomposition of $G$ into
its irreducible components. If $G_j \subset F$, then we can remove
it from the list. Thus we may assume $G_j$ is not contained in $F$
for $j = 1, \ldots, m$.
For $j = 1, \ldots, m$ choose $y_j \in G_j$ with $y_j \not \in F$
and $y_j \not \in G_{j'}$ for $j' \not = j$. This is possible as
there are no inclusions among the irreducible components of $G$.
Choose $w_j \in V(y_j)$ not contained in the span of the images of
$s_1(y_j), \ldots, s_h(y_j)$; this is possible because
$h + k \leq \dim V(y_j)$ and $k > 0$.

\medskip\noindent
Apply the induction hypothesis to the $h + 1$ sections
$s_1, \ldots, s_h, u$, the closed set $F \cup G$,
the points $x_1, \ldots, x_n, y_1, \ldots, y_m \in F \cup G$,
the elements $0 \in V(x_i)$ and $w_j \in V(y_j)$, and
the integer $k - 1$. Note that we have increased $h$ by $1$ and
decreased $k$ by $1$ hence the assumption $(*)$ of the proposition
remains valid. This produces $t \in M$ and $H \subset \Omega$ closed
satisfying (a), (b), (c) for $s_1, \ldots, s_h, u$, $F \cup G$,
$x_1, \ldots, x_n, y_1, \ldots, y_m$, $0, \ldots, 0, w_1, \ldots, w_m$,
and $k - 1$.

\medskip\noindent
Let $H_1, \ldots, H_p \subset H$ be the irreducible components of $H$
which are not contained in $F \cup G$. As before pick $z_l \in H_l$,
$z_l \not \in F \cup G$ and $z_l \not \in H_{l'}$ for $l' \not = l$.
Using Algebra, Lemma \href{https://stacks.math.columbia.edu/tag/00DT}{lemma chinese remainder (Tag 00DT)}
we may choose $f \in R$ such that
$f(y_j) = 1$, $j = 1, \ldots, m$ and $f(z_l) = 0$, $l = 1, \ldots, p$.
Claim: the element $s = u + f t$ works.

\medskip\noindent
First, the value $s(x_i)$ agrees with $u(x_i)$ because
$t(x_i) = 0$ and hence we see that $s(x_i)$ maps to $v_i$.
This proves (a).
To finish the proof it suffices to show that every irreducible
component $Z$ of $Z(s_1, \ldots, s_h, s)$ not contained in $F$
has codimension $\geq k$ in $\Omega$. Namely, then we can set $F'$
equal to the union of these and we get (b) and (c).
We can see that irreducible components $Z$ of $Z(s_1, \ldots, s_h, s)$
of codimension $\leq k - 1$ do not exist as follows:
\begin{enumerate}
\item Observe that
$Z(s_1, \ldots, s_h, s) \subset Z(s_1, \ldots, s_h, u, t) = F \cup H$
as $s = u + ft$. Hence $Z \subset H$.
\item The irreducible components of $H$ have codimension $\geq k - 1$.
Hence $Z$ is equal to an irreducible component of $H$
as $Z$ has codimension $\leq k - 1$.
Hence $Z = H_l$ for some $l \in \{1, \ldots, p\}$ or
$Z = G_j$ for some $j \in \{1, \ldots m\}$.
\item But $Z = G_j$ is impossible as $s_1(y_j), \ldots, s_h(y_j)$
map to linearly independent elements of $V(y_j)$ and
$s(y_j) = u(y_j) + f(y_j) t(y_j) = u(y_j) + t(y_j)$
maps to an element of the form
$$
\text{linear combination images of }s_i(y_j) + w_j
$$
which is linearly independent of the images of $s_1(y_j), \ldots, s_h(y_j)$
in $V(y_j)$ by our choice of $w_j$.
\item Also $Z = H_l$ is impossible. Namely, again
$s_1(z_l), \ldots, s_h(z_l)$ map to linearly independent
elements of $V(z_l)$ and $s(z_l) = u(z_l) + f(z_l) t(z_l) = u(z_l)$
maps to an element of $V(z_l)$ linearly independent of those
as $z_l \not \in F \cup G$.
\end{enumerate}
This finishes the proof.
\end{proof}
\end{document}
